\section*{Orientaciones bibliográficas}
\addcontentsline{toc}{section}{Orientaciones bibliográficas}
Para extensiones simples, polinomios mínimos y automorfismos, véase \citet[cap.~VII, apartados 1.1--2, definiciones 1.1--2, proposiciones 1.3/5, ejemplo 1.4 y corolario 1.7]{aluffi2009}. La construcción de una extensión que contiene una raíz se trata en V.5.7.

Para extensiones finitas y algebraicas, la fórmula de la torre y los ejemplos de extensiones algebraicas, véase \citet[cap.~VII, apartado 1.3, definición 1.8, lema 1.9 y resultados 1.10--19]{aluffi2009}. Las comprobaciones de expresiones racionales y de automorfismos sobre raíces corresponden a los ejercicios VII.1.3 y VII.1.6. La construcción y el conteo de automorfismos simples se desarrollan en la sección~\ref{sec:campos-simples}.

Para clausuras algebraicas, véase \citet[cap.~VII, apartado 2.1, lema 2.1, definición 2.2, teorema 2.3 y resultados 2.4--6/8]{aluffi2009}; la comprobación sobre sus raíces corresponde al ejercicio VII.2.2. Los campos de descomposición y la normalidad se tratan en VII.4.1--4/7--8. El diagrama de comparación de raíces se desarrolla en la sección~\ref{tema:25-1-2}.

Para normalidad y el diagrama de comparación, véase \citet[cap.~VII, definición 4.7 y teorema 4.8]{aluffi2009}. La separabilidad, Frobenius y los campos perfectos se desarrollan en VII.4.11--15/17--21. Las comprobaciones de la identidad binomial, la independencia del campo de descomposición y la recíproca sobre campos perfectos corresponden a los ejercicios VII.4.8, VII.4.1 y VII.4.12, respectivamente. Eisenstein y Gauss se aplican en las formas de las secciones~\ref{sec:anillos-fracciones} y~\ref{sec:anillos-irreducibilidad}.

El grado separable se estudia en \citet[cap.~VII, lemas 4.22--23 y proposición 4.24]{aluffi2009}. Para el elemento primitivo, los campos intermedios y el conteo de automorfismos, véanse VII.5.18--20. La comprobación para campos finitos se desarrolla en el \cref{ejer:C26-campo-finito-simple}. Los campos fijos y las propiedades iniciales de la correspondencia se encuentran en VII.6.1--4; el último ejemplo corresponde a VII.6.4.

En \citet[cap.~VII]{aluffi2009}, los campos fijos y las extensiones de Galois se estudian en 6.5--6/9--11; los retículos y grupos cociente, en 6.12/15; la acción en raíces y los discriminantes, en 7.17--19. Los ejercicios resueltos corresponden a VII.4.7 y VII.6.3.

La adjunción de raíces y las caracterizaciones de los campos algebraicamente cerrados se desarrollan en \citet[cap.~V, apartados 5.1--4 y resultados 5.7--10]{aluffi2009}.
